
      \documentclass[fleqn]{article}      
      \usepackage{amsmath}
      \usepackage{fontspec}
      \usepackage{kotex} % 한국어 지원  

      \begin{document}      
       객관식 질문

쉬움:
1. 행렬의 크기를 정할 때, 행(row)의 개수와 열(column)의 개수는 무엇을 의미하나요?  
   a) 행렬의 원소 개수  
   b) 행렬의 차원  
   c) 행렬의 합  

2. 다음 중 행렬의 덧셈이 가능한 경우는?  
   a) 2x3 행렬과 2x2 행렬  
   b) 3x2 행렬과 3x2 행렬  
   c) 2x1 행렬과 1x2 행렬  

3. 행렬 A의 전치 행렬을 무엇이라고 하나요?  
   a) A’  
   b) A^T  
   c) A^-1  

중간:
1. 다음 중 행렬 곱셈에 대한 설명으로 옳은 것은?  
   a) 두 행렬의 크기가 같아야 곱할 수 있다.  
   b) 첫 번째 행렬의 열 수와 두 번째 행렬의 행 수가 같아야 곱할 수 있다.  
   c) 행렬 곱셈은 항상 교환 법칙이 성립한다.  

2. 다음 중 단위 행렬의 성질은 무엇인가요?  
   a) 모든 원소가 0이다.  
   b) 대각선 원소가 1이고 나머지는 0이다.  
   c) 모든 원소가 1이다.  

3. 두 행렬 A와 B가 같으려면 어떤 조건이 필요한가요?  
   a) 크기가 같고 모든 원소가 같아야 한다.  
   b) 크기가 달라도 된다.  
   c) 단지 행렬의 크기만 같으면 된다.  

어려움:
1. 행렬 A = [[1, 2], [3, 4]]와 B = [[5, 6], [7, 8]]의 곱 AB는 무엇인가요?  
   a) [[19, 22], [43, 50]]  
   b) [[23, 34], [31, 46]]  
   c) [[43, 50], [19, 22]]  

2. 주어진 행렬 A = [[1, 0, 2], [0, 1, 3]]의 전치 행렬 A^T를 구하시오.  
   a) [[1, 0], [0, 1], [2, 3]]  
   b) [[1, 0], [2, 3]]  
   c) [[1, 0], [0, 1], [2, 0]]  

3. 행렬 A의 역행렬이 존재하기 위한 조건은 무엇인가요?  
   a) A의 행렬식이 0이 아니다.  
   b) A의 모든 원소가 0이 아니다.  
   c) A의 행과 열이 같아야 한다.  

 주관식 질문

쉬움:  
1. 행렬이란 무엇인가요?  
2. 행렬의 덧셈을 어떻게 수행하나요?  
3. 행렬의 전치란 무엇인가요?  

중간:  
1. 2x3 행렬과 3x2 행렬의 곱을 설명하세요.  
2. 단위 행렬의 정의와 그 성질을 설명하세요.  
3. 행렬의 크기를 정의하세요.  
   
어려움:  
1. 두 행렬 A와 B의 곱을 구하는 방법을 설명하세요.  
2. 행렬의 역행렬을 구할 수 있는 경우와 그 방법에 대해 설명하세요.  
3. 행렬식을 계산하는 방법을 설명하세요.  

      
      \end{document} 
  